% questions.tex —— 2019 年·全国Ⅰ卷·理科 · 题目内容片段（无答案）
% 由 tools/generate_projects.py 依来源总账自 2019/2019_数学真题_理_试题.md 生成，勿手改（改源 MD 后重跑）


\begin{gknotice}
\textbf{ 卷种提示}：本卷为文理分科时代的旧结构，含\textbf{现已取消}的选做题（第 22/23 题）与\textbf{已删除}的程序框图考点（第 8 题）。适合\textbf{按知识点抽题}，不适合整套限时模拟。

\textbf{适用省份}：河南、河北、山西、江西、湖北、湖南、广东、安徽、福建、山东
\end{gknotice}

\par\medskip

\gkpart{一、选择题}

\begin{gktip}
本题共 12 小题，每小题 5 分，共 60 分。
\end{gktip}

\begin{question}
\gkti{1}已知集合 $M=\{x\mid-4<x<2\}$，$N=\left\{x\mid x^{2}-x-6<0\right\}$，则 $M\cap N=$\mbox{（　　）}

\begin{gkopts}(2)
\task A．$\{x\mid-4<x<3\}$
\task B．$\{x\mid-4<x<-2\}$
\task C．$\{x\mid-2<x<2\}$
\task D．$\{x\mid2<x<3\}$
\end{gkopts}

\end{question}
\begin{question}
\gkti{2}设复数 $z$ 满足 $|z-\mathrm i|=1$，$z$ 在复平面内对应的点为 $(x,y)$，则\mbox{（　　）}

\begin{gkopts}(2)
\task A．$(x+1)^{2}+y^{2}=1$
\task B．$(x-1)^{2}+y^{2}=1$
\task C．$x^{2}+(y-1)^{2}=1$
\task D．$x^{2}+(y+1)^{2}=1$
\end{gkopts}

\end{question}
\begin{question}
\gkti{3}已知 $a=\log_2 0.2$，$b=2^{0.2}$，$c=0.2^{0.3}$，则\mbox{（　　）}

\begin{gkopts}(4)
\task A．$a<b<c$
\task B．$a<c<b$
\task C．$c<a<b$
\task D．$b<c<a$
\end{gkopts}

\end{question}
\begin{question}
\gkti{4}古希腊时期，人们认为最美人体的头顶至肚脐的长度与肚脐至足底的长度之比是 $\dfrac{\sqrt5-1}2\ (\approx0.618$，称为黄金分割比例$)$，著名的"断臂维纳斯"便是如此。此外，最美人体的头顶至咽喉的长度与咽喉至肚脐的长度之比也是 $\dfrac{\sqrt5-1}2$。若某人满足上述两个黄金分割比例，且腿长为 $105\ \mathrm{cm}$，头顶至脖子下端的长度为 $26\ \mathrm{cm}$，则其身高可能是\mbox{（　　）}

\begin{gkopts}(2)
\task A．$165\ \mathrm{cm}$
\task B．$175\ \mathrm{cm}$
\task C．$185\ \mathrm{cm}$
\task D．$190\ \mathrm{cm}$
\end{gkopts}

\Needspace{0.4\textheight}
\end{question}
\begin{question}
\gkti{5}函数 $f(x)=\dfrac{\sin x+x}{\cos x+x^{2}}$ 在 $[-\pi,\pi]$ 的图象大致为\mbox{（　　）}
\nopagebreak[4]
\begin{center}% 2019 全国Ⅰ卷 理5/文5：f(x)=(sin x+x)/(cos x+x^2) 图象四选一
% (C)(D) 由精确表达式绘制（D=f 为正确图象，C=-f）；
% (A)(B) 为干扰项，依据原卷形态用样条示意重绘（原卷亦为定性示意图）
\begin{tikzpicture}[font=\footnotesize]
\pgfplotsset{
  panel/.style={width=4.6cm, height=3.4cm, axis lines=middle,
    xmin=-3.6, xmax=3.6, ymin=-1.6, ymax=1.7,
    xtick={-3.14159,3.14159}, xticklabels={$-\pi$,$\pi$},
    ytick={1}, yticklabels={$1$},
    xlabel={$x$}, ylabel={$y$},
    every axis x label/.style={at={(current axis.right of origin)}, anchor=west},
    every axis y label/.style={at={(current axis.above origin)}, anchor=south}}}
  \begin{axis}[panel, at={(0,0)}, name=axA]
    \addplot[smooth, thick] coordinates
      {(-3.14,0.55) (-2.2,0.42) (-1.2,0.18) (-0.5,-0.15) (0,0)
       (0.45,0.75) (0.8,1.05) (1.6,0.8) (2.4,0.4) (3.14,0.1)};
    \node[below left] at (axis cs:0,0) {$O$};
  \end{axis}
  \node[below=0.5mm] at (axA.south) {(A)};
  \begin{axis}[panel, at={(5.5cm,0)}, name=axB]
    \addplot[smooth, thick] coordinates
      {(-3.14,-0.02) (-2.4,-0.25) (-1.5,-0.45) (-0.6,-0.3) (0,0)
       (0.7,0.45) (1.3,0.58) (2.2,0.35) (3.14,0.12)};
    \node[below left] at (axis cs:0,0) {$O$};
  \end{axis}
  \node[below=0.5mm] at (axB.south) {(B)};
  \begin{axis}[panel, at={(0,-4.4cm)}, name=axC]
    \addplot[domain=-3.14159:3.14159, samples=200, thick]
      {-(sin(deg(x))+x)/(cos(deg(x))+x^2)};
    \node[below left] at (axis cs:0,0) {$O$};
  \end{axis}
  \node[below=0.5mm] at (axC.south) {(C)};
  \begin{axis}[panel, at={(5.5cm,-4.4cm)}, name=axD]
    \addplot[domain=-3.14159:3.14159, samples=200, thick]
      {(sin(deg(x))+x)/(cos(deg(x))+x^2)};
    \node[below left] at (axis cs:0,0) {$O$};
  \end{axis}
  \node[below=0.5mm] at (axD.south) {(D)};
\end{tikzpicture}
\end{center}
\nopagebreak[4]
\end{question}
\begin{question}
\gkti{6}我国古代典籍《周易》用"卦"描述万物的变化。每一"重卦"由从下到上排列的 $6$ 个爻组成，爻分为阳爻和阴爻。在所有重卦中随机取一重卦，则该重卦恰有 $3$ 个阳爻的概率是\mbox{（　　）}

\begin{gkopts}(4)
\task A．$\dfrac5{16}$
\task B．$\dfrac{11}{32}$
\task C．$\dfrac{21}{32}$
\task D．$\dfrac{11}{16}$
\end{gkopts}

\end{question}
\begin{question}
\gkti{7}已知非零向量 $\boldsymbol a$，$\boldsymbol b$ 满足 $|\boldsymbol a|=2|\boldsymbol b|$，且 $(\boldsymbol a-\boldsymbol b)\perp\boldsymbol b$，则 $\boldsymbol a$ 与 $\boldsymbol b$ 的夹角为\mbox{（　　）}

\begin{gkopts}(4)
\task A．$\dfrac\pi6$
\task B．$\dfrac\pi3$
\task C．$\dfrac{2\pi}3$
\task D．$\dfrac{5\pi}6$
\end{gkopts}

\Needspace{0.4\textheight}
\end{question}
\begin{question}
\gkti{8}如图是求 $\dfrac1{2+\dfrac1{2+\dfrac12}}$ 的程序框图，图中空白框中应填入\mbox{（　　）}
\nopagebreak[4]
\begin{center}% 2019 全国Ⅰ卷 理8/文9：求连分数 1/(2+1/(2+1/2)) 的程序框图
% 空白矩形为题目要求填入的位置（正确答案 A=1/(2+A)）
\begin{tikzpicture}[font=\small, >=Stealth, node distance=7mm,
  box/.style={draw, minimum height=1.9em, inner xsep=7pt},
  term/.style={draw, rounded corners=7pt, minimum height=1.9em, inner xsep=9pt},
  io/.style={draw, trapezium, trapezium left angle=70, trapezium right angle=110,
             minimum height=1.9em, inner xsep=2pt},
  dec/.style={draw, diamond, aspect=2.4, inner sep=1.5pt}]
  \node[term] (start) {开始};
  \node[box, below=of start] (init) {$A=\dfrac12$};
  \node[box, below=of init] (k1) {$k=1$};
  \coordinate[below=4mm of k1] (merge);
  \node[dec, below=2mm of merge] (cond) {$k\le2$};
  \node[box, below=of cond, minimum width=5.6em] (blank) {\phantom{$A=\dfrac1{2+A}$}};
  \node[box, below=of blank] (inc) {$k=k+1$};
  \node[io, right=14mm of cond] (out) {输出 $A$};
  \node[term, below=of out] (end) {结束};
  \draw[->] (start) -- (init);
  \draw[->] (init) -- (k1);
  \draw[->] (k1) -- (cond);
  \draw[->] (cond) -- node[left]{是} (blank);
  \draw[->] (blank) -- (inc);
  \draw[->] (inc.west) -- ++(-8mm,0) |- (merge);
  \draw[->] (cond) -- node[above]{否} (out);
  \draw[->] (out) -- (end);
\end{tikzpicture}
\end{center}
\nopagebreak[4]
\begin{gkopts}(2)
\task A．$A=\dfrac1{2+A}$
\task B．$A=2+\dfrac1A$
\task C．$A=\dfrac1{1+2A}$
\task D．$A=1+\dfrac1{2A}$
\end{gkopts}

\begin{gktip}
\textbf{注}：程序框图（算法初步）在现行新课标中已删除。
\end{gktip}

\end{question}
\begin{question}
\gkti{9}记 $S_n$ 为等差数列 $\left\{a_n\right\}$ 的前 $n$ 项和。已知 $S_4=0$，$a_5=5$，则\mbox{（　　）}

\begin{gkopts}(2)
\task A．$a_n=2n-5$
\task B．$a_n=3n-10$
\task C．$S_n=2n^{2}-8n$
\task D．$S_n=\dfrac12n^{2}-2n$
\end{gkopts}

\end{question}
\begin{question}
\gkti{10}已知椭圆 $C$ 的焦点为 $F_1(-1,0)$，$F_2(1,0)$，过 $F_2$ 的直线与 $C$ 交于 $A$，$B$ 两点。若 $\left|AF_2\right|=2\left|F_2B\right|$，$|AB|=\left|BF_1\right|$，则 $C$ 的方程为\mbox{（　　）}

\begin{gkopts}(2)
\task A．$\dfrac{x^{2}}2+y^{2}=1$
\task B．$\dfrac{x^{2}}3+\dfrac{y^{2}}2=1$
\task C．$\dfrac{x^{2}}4+\dfrac{y^{2}}3=1$
\task D．$\dfrac{x^{2}}5+\dfrac{y^{2}}4=1$
\end{gkopts}

\end{question}
\begin{question}
\gkti{11}关于函数 $f(x)=\sin|x|+|\sin x|$ 有下述四个结论：

  ① $f(x)$ 是偶函数；
  ② $f(x)$ 在区间 $\left(\dfrac\pi2,\pi\right)$ 单调递增；
  ③ $f(x)$ 在 $[-\pi,\pi]$ 有 $4$ 个零点；
  ④ $f(x)$ 的最大值为 $2$。

其中所有正确结论的编号是\mbox{（　　）}

\begin{gkopts}(4)
\task A．①②④
\task B．②④
\task C．①④
\task D．①③
\end{gkopts}

\end{question}
\begin{question}
\gkti{12}已知三棱锥 $P\text{-}ABC$ 的四个顶点在球 $O$ 的球面上，$PA=PB=PC$，$\triangle ABC$ 是边长为 $2$ 的正三角形，$E$，$F$ 分别是 $PA$，$AB$ 的中点，$\angle CEF=90^\circ$，则球 $O$ 的体积为\mbox{（　　）}

\begin{gkopts}(4)
\task A．$8\sqrt6\pi$
\task B．$4\sqrt6\pi$
\task C．$2\sqrt6\pi$
\task D．$\sqrt6\pi$
\end{gkopts}

\par\medskip

\end{question}
\gkpart{二、填空题}

\begin{gktip}
本题共 4 小题，每小题 5 分，共 20 分。
\end{gktip}

\begin{question}
\gkti{13}曲线 $y=3\left(x^{2}+x\right)\mathrm e^{x}$ 在点 $(0,0)$ 处的切线方程为 \gkblank{}．

\end{question}
\begin{question}
\gkti{14}记 $S_n$ 为等比数列 $\left\{a_n\right\}$ 的前 $n$ 项和。若 $a_1=\dfrac13$，$a_4^{2}=a_6$，则 $S_5=$ \gkblank{}．

\end{question}
\begin{question}
\gkti{15}甲、乙两队进行篮球决赛，采取七场四胜制（当一队赢得四场胜利时，该队获胜，决赛结束）。根据前期比赛成绩，甲队的主客场安排依次为"主主客客主客主"。设甲队主场取胜的概率为 $0.6$，客场取胜的概率为 $0.5$，且各场比赛结果相互独立，则甲队以 $4:1$ 获胜的概率是 \gkblank{}．

\end{question}
\begin{question}
\gkti{16}已知双曲线 $C:\dfrac{x^{2}}{a^{2}}-\dfrac{y^{2}}{b^{2}}=1\ (a>0,\ b>0)$ 的左、右焦点分别为 $F_1$，$F_2$，过 $F_1$ 的直线与 $C$ 的两条渐近线分别交于 $A$，$B$ 两点。若 $\overrightarrow{F_1A}=\overrightarrow{AB}$，$\overrightarrow{F_1B}\cdot\overrightarrow{F_2B}=0$，则 $C$ 的离心率为 \gkblank{}．

\par\medskip

\end{question}
\gkpart{三、解答题}

\begin{gktip}
共 70 分。第 17～21 题为必考题，第 22、23 题为选考题。
\end{gktip}

\gksubpart{（一）必考题：共 60 分}

\begin{question}
\gkti{17}\textbf{（12 分）} $\triangle ABC$ 的内角 $A$，$B$，$C$ 的对边分别为 $a$，$b$，$c$，设 $(\sin B-\sin C)^{2}=\sin^{2}A-\sin B\sin C$。

  (1) 求 $A$；

  (2) 若 $\sqrt2a+b=2c$，求 $\sin C$。

\Needspace{0.4\textheight}
\end{question}
\begin{question}
\gkti{18}\textbf{（12 分）} 如图，直四棱柱 $ABCD\text{-}A_1B_1C_1D_1$ 的底面是菱形，$AA_1=4$，$AB=2$，$\angle BAD=60^\circ$，$E$，$M$，$N$ 分别是 $BC$，$BB_1$，$A_1D$ 的中点。
\nopagebreak[4]
\begin{center}% 2019 全国Ⅰ卷 理18/文19：直四棱柱 ABCD-A1B1C1D1（底面菱形），
% E、M、N 分别为 BC、BB1、A1D 的中点
\begin{tikzpicture}[scale=0.92, font=\small]
  % 底面菱形（斜二测示意）：A 前左、B 前右、C 右后、D 左后
  \coordinate (A) at (0,0);
  \coordinate (B) at (2.1,-0.25);
  \coordinate (C) at (3.15,0.6);
  \coordinate (D) at (1.05,0.85);
  \coordinate (A1) at (0,3.0);
  \coordinate (B1) at (2.1,2.75);
  \coordinate (C1) at (3.15,3.6);
  \coordinate (D1) at (1.05,3.85);
  \coordinate (E) at ($(B)!0.5!(C)$);
  \coordinate (M) at ($(B)!0.5!(B1)$);
  \coordinate (N) at ($(A1)!0.5!(D)$);
  \draw (A) -- (B) -- (C);
  \draw[dashed] (C) -- (D) -- (A);
  \draw (A) -- (A1)  (B) -- (B1)  (C) -- (C1);
  \draw[dashed] (D) -- (D1);
  \draw (A1) -- (B1) -- (C1) -- (D1) -- cycle;
  \draw[dashed] (A1) -- (D) ;
  \draw[dashed] (N) -- (M);
  \draw[dashed] (D) -- (E);
  \draw (C1) -- (E);
  \draw[dashed] (C1) -- (D);
  \node[below left] at (A) {$A$};
  \node[below] at (B) {$B$};
  \node[right] at (C) {$C$};
  \node[below right, inner sep=0.5pt] at (D) {$D$};
  \node[above left, inner sep=1pt] at (A1) {$A_1$};
  \node[below right, inner sep=1pt] at (B1) {$B_1$};
  \node[right] at (C1) {$C_1$};
  \node[above] at (D1) {$D_1$};
  \node[below right, inner sep=1pt] at (E) {$E$};
  \node[right, inner sep=2pt] at (M) {$M$};
  \node[left, inner sep=2.5pt] at (N) {$N$};
\end{tikzpicture}
\end{center}
\nopagebreak[4]
  (1) 证明：$MN\parallel$ 平面 $C_1DE$；

  (2) 求二面角 $A\text{-}MA_1\text{-}N$ 的正弦值。

\end{question}
\begin{question}
\gkti{19}\textbf{（12 分）} 已知抛物线 $C:y^{2}=3x$ 的焦点为 $F$，斜率为 $\dfrac32$ 的直线 $l$ 与 $C$ 的交点为 $A$，$B$，与 $x$ 轴的交点为 $P$。

  (1) 若 $|AF|+|BF|=4$，求 $l$ 的方程；

  (2) 若 $\overrightarrow{AP}=3\overrightarrow{PB}$，求 $|AB|$。

\end{question}
\begin{question}
\gkti{20}\textbf{（12 分）} 已知函数 $f(x)=\sin x-\ln(1+x)$，$f'(x)$ 为 $f(x)$ 的导数。证明：

  (1) $f'(x)$ 在区间 $\left(-1,\dfrac\pi2\right)$ 存在唯一极大值点；

  (2) $f(x)$ 有且仅有 $2$ 个零点。

\end{question}
\begin{question}
\gkti{21}\textbf{（12 分）} 为治疗某种疾病，研制了甲、乙两种新药，希望知道哪种新药更有效，为此进行动物试验。试验方案如下：每一轮选取两只白鼠对药效进行对比试验。对于两只白鼠，随机选一只施以甲药，另一只施以乙药。一轮的治疗结果得出后，再安排下一轮试验。当其中一种药治愈的白鼠比另一种药治愈的白鼠多 $4$ 只时，就停止试验，并认为治愈只数多的药更有效。为了方便描述问题，约定：对于每轮试验，若施甲药的白鼠治愈且施乙药的白鼠未治愈则甲药得 $1$ 分，乙药得 $-1$ 分；若施乙药的白鼠治愈且施甲药的白鼠未治愈则乙药得 $1$ 分，甲药得 $-1$ 分；若都治愈或都未治愈则两种药均得 $0$ 分。甲、乙两种药的治愈率分别记为 $\alpha$ 和 $\beta$，一轮试验中甲药的得分记为 $X$。

  (1) 求 $X$ 的分布列；

  (2) 若甲药、乙药在试验开始时都赋予 $4$ 分，$p_i\ (i=0,1,\cdots,8)$ 表示"甲药的累计得分为 $i$ 时，最终认为甲药比乙药更有效"的概率，则 $p_0=0$，$p_8=1$，$p_i=ap_{i-1}+bp_i+cp_{i+1}\ (i=1,2,\cdots,7)$，其中 $a=P(X=-1)$，$b=P(X=0)$，$c=P(X=1)$。假设 $\alpha=0.5$，$\beta=0.8$。

  　① 证明：$\left\{p_{i+1}-p_i\right\}\ (i=0,1,2,\cdots,7)$ 为等比数列；

  　② 求 $p_4$，并根据 $p_4$ 的值解释这种试验方案的合理性。

\par\medskip

\end{question}
\gksubpart{（二）选考题：共 10 分，第 22、23 题任选一题作答}

\begin{gktip}
\textbf{ 现行考纲已取消选做题}，以下两题仅作历史资料保留。
\end{gktip}

\begin{question}
\gkti{22}\textbf{（10 分）〔选修 4-4：坐标系与参数方程〕} （\textbf{已取消}）

在直角坐标系 $xOy$ 中，曲线 $C$ 的参数方程为 $\begin{cases}x=\dfrac{1-t^{2}}{1+t^{2}}\\[2mm] y=\dfrac{4t}{1+t^{2}}\end{cases}$（$t$ 为参数）。以坐标原点 $O$ 为极点，$x$ 轴的正半轴为极轴建立极坐标系，直线 $l$ 的极坐标方程为 $2\rho\cos\theta+\sqrt3\rho\sin\theta+11=0$。

  (1) 求 $C$ 和 $l$ 的直角坐标方程；

  (2) 求 $C$ 上的点到 $l$ 距离的最小值。

\end{question}
\begin{question}
\gkti{23}\textbf{（10 分）〔选修 4-5：不等式选讲〕} （\textbf{已取消}）

已知 $a$，$b$，$c$ 为正数，且满足 $abc=1$。证明：

  (1) $\dfrac1a+\dfrac1b+\dfrac1c\le a^{2}+b^{2}+c^{2}$；

  (2) $(a+b)^{3}+(b+c)^{3}+(c+a)^{3}\ge24$。

\par\medskip


\end{question}